\section{Risultati preliminari}

Nel primo esperimento intermedio, la rete ha ottenuto le seguenti metriche sul test set:

\begin{table}[H]
    \centering
    \begin{tabular}{lr}
        \toprule
        Metrica & Valore \\
        \midrule
        MAE & 0.6952 \\
        RMSE & 0.9114 \\
        $R^2$ & 0.9985 \\
        MAPE, prezzi $> 1$ & 10.0137\% \\
        \bottomrule
    \end{tabular}
    \caption{Metriche preliminari della rete neurale rispetto ai prezzi Black-Scholes.}
\end{table}

Il valore elevato di $R^2$ suggerisce che la rete neurale riesce gia' ad approssimare molto bene la funzione di pricing Black-Scholes.
Gli errori assoluti medi rimangono contenuti rispetto al range complessivo dei prezzi.
Il MAPE e' piu' delicato da interpretare, perche' i prezzi di opzioni molto out-of-the-money possono essere vicini a zero; per questo viene calcolato solo su prezzi maggiori di 1.

\begin{figure}[H]
    \centering
    \maybeincludegraphics[width=0.72\linewidth]{../outputs/figures/loss.png}
    \caption{Training loss e validation loss.}
\end{figure}

\begin{figure}[H]
    \centering
    \maybeincludegraphics[width=0.62\linewidth]{../outputs/figures/true_vs_predicted.png}
    \caption{Prezzi Black-Scholes rispetto ai prezzi predetti dalla rete.}
\end{figure}

\begin{figure}[H]
    \centering
    \maybeincludegraphics[width=0.72\linewidth]{../outputs/figures/error_distribution.png}
    \caption{Distribuzione degli errori di predizione.}
\end{figure}

\begin{figure}[H]
    \centering
    \begin{subfigure}{0.32\linewidth}
        \centering
        \maybeincludegraphics[width=\linewidth]{../outputs/figures/error_vs_moneyness.png}
        \caption{Moneyness}
    \end{subfigure}
    \begin{subfigure}{0.32\linewidth}
        \centering
        \maybeincludegraphics[width=\linewidth]{../outputs/figures/error_vs_maturity.png}
        \caption{Maturity}
    \end{subfigure}
    \begin{subfigure}{0.32\linewidth}
        \centering
        \maybeincludegraphics[width=\linewidth]{../outputs/figures/error_vs_volatility.png}
        \caption{Volatilita'}
    \end{subfigure}
    \caption{Errore assoluto rispetto ad alcune feature finanziarie.}
\end{figure}

\subsection{Confronto Monte Carlo}

Sul sottoinsieme usato per il benchmark Monte Carlo, la simulazione numerica ha ottenuto:

\begin{table}[H]
    \centering
    \begin{tabular}{lr}
        \toprule
        Metrica & Valore \\
        \midrule
        MAE & 0.3100 \\
        RMSE & 0.5201 \\
        $R^2$ & 0.9996 \\
        MAPE, prezzi $> 1$ & 2.2139\% \\
        \bottomrule
    \end{tabular}
    \caption{Metriche preliminari Monte Carlo rispetto ai prezzi Black-Scholes.}
\end{table}

Il confronto con Monte Carlo e' utile per mostrare la differenza tra un metodo numerico basato su simulazione e un modello supervisionato che, una volta addestrato, produce predizioni molto rapidamente.

